
\begin{figure}[!htbp]
\centering
\begin{subfigure}[t]{0.48\textwidth}\centering
\includegraphics[width=\linewidth]{figures/ch4/cp112_generated/gii_max_drift.png}\subcaption{ULS drift}
\end{subfigure}\hfill
\begin{subfigure}[t]{0.48\textwidth}\centering
\includegraphics[width=\linewidth]{figures/ch4/cp112_generated/gii_cmcr.png}\subcaption{CM--CR separation}
\end{subfigure}

\begin{subfigure}[t]{0.48\textwidth}\centering
\includegraphics[width=\linewidth]{figures/ch4/cp112_generated/gii_peak_pushover.png}\subcaption{Peak pushover shear}
\end{subfigure}\hfill
\begin{subfigure}[t]{0.48\textwidth}\centering
\includegraphics[width=\linewidth]{figures/ch4/cp112_generated/gii_pp_displacement.png}\subcaption{Performance-point roof displacement}
\end{subfigure}
\caption{Geometric Irregularity Index (GII) and selected seismic responses.}
\label{fig:gii_multiresponse_cp112a}
\end{figure}

Figure~\ref{fig:gii_multiresponse_cp112a} demonstrates that the same or similar GII can correspond to different drift, CM--CR separation, peak pushover strength and performance-point displacement. The plotted responses therefore do not support treating GII as a unique predictor of seismic behaviour.


The performance-point-shear comparison is already presented in Figure~\ref{fig:gii_pp_base_shear}. Together with the four-response comparison in Figure~\ref{fig:gii_multiresponse_cp112a}, these results indicate that plan topology and stiffness distribution cannot be described by area deficit alone.

\FloatBarrier

\begin{table}[!htbp]
\centering
\caption{Final seven-model comparative summary.}
\label{tab:final_summary_cp112}
\scriptsize
\resizebox{\textwidth}{!}{%
\begin{tabular}{@{}lrrrrrrp{2.7cm}@{}}
\toprule
Model & GII & Max drift & Max CM--CR & EQx shear & Peak push. & Max PP disp. & Hinge evidence \\
 & (\%) & ratio & (m) & (kN) & (kN) & (mm) & \\
\midrule
Regular & 0.00 & 0.004784 & 0.398 & 532.489 & 1836.338 & 26.053 & Final-step X/Y images; member-level table unavailable \\
L & 16.67 & 0.004758 & 0.581 & 476.287 & 1800.120 & 25.900 & Same evidence boundary \\
L1 & 33.33 & 0.004570 & 0.640 & 403.859 & 1606.250 & 24.592 & Same evidence boundary \\
L2 & 33.33 & 0.004025 & 1.186 & 456.660 & 1804.152 & 21.171 & Same evidence boundary \\
T & 33.33 & 0.004618 & 0.475 & 407.396 & 1640.715 & 24.802 & Same evidence boundary \\
U & 16.67 & 0.004954 & 0.424 & 504.048 & 1882.070 & 27.072 & Same evidence boundary \\
Non-Parallel & 22.22 & 0.004041 & 1.237 & 486.757 & 1827.694 & 21.328 & Same evidence boundary \\
\bottomrule
\end{tabular}
}
\par\smallskip
\begin{minipage}{0.97\textwidth}\footnotesize
EQx shear is the validated Soil C, $Z=0.40$ Story-1 value. Formal NBC torsional-irregularity classification is not assigned from CM--CR separation. The table summarizes the evidence used in this thesis; it is not a code-acceptance table.
\end{minipage}
\end{table}

\clearpage
\subsection{Alignment of Research Questions, Results and Conclusions}
The six research objectives are linked to their corresponding Chapter~4 evidence and Chapter~5 findings in Table~\ref{tab:rq_objective_traceability_cp112}. This mapping provides a guide to the manuscript argument; it does not constitute an additional dataset.
\begin{table}[!ht]
\centering
\caption{Alignment of the six research questions/objectives with results and conclusions.}
\label{tab:rq_objective_traceability_cp112}
\scriptsize
\begin{tabular}{@{}c p{8.0cm} p{3.4cm}@{}}
\toprule
No. & Result evidence in Chapter 4 & Conclusion linkage in Chapter 5 \\
\midrule
1 & Phase I concrete design-check evidence for the 2012 reference & 2012 design-check finding \\
2 & Regular benchmark modal/storey/drift/pushover/performance response & Regular benchmark finding \\
3 & Seven-model plan-irregular comparison and complete drift/force/capacity tables & Irregular-model finding \\
4 & CM--CR coordinates, drift profiles, pushover and performance-point comparisons & CM--CR/response finding \\
5 & Phase IV applicability-boundary synthesis and final comparative table & RUD scope/height finding \\
6 & Seven-model GII tables and selected GII--response comparisons & GII response finding \\
\bottomrule
\end{tabular}
\end{table}
\FloatBarrier
